\chapter[Resultados]{Resultados}
\label{chap: resultados}

Este capítulo apresenta os resultados da bateria experimental descrita no \autoref{chap: protocolo}. A \autoref{sec: res_partida} estabelece os pontos de referência: o elenco canônico, do qual a busca parte, e os modelos nulos, que dão escala às métricas. As três seções seguintes respondem, cada uma, a uma questão de pesquisa: o equilíbrio (\autoref{sec: res_equilibrio}), a identidade (\autoref{sec: res_identidade}) e o compromisso entre os dois (\autoref{sec: res_compromisso}). A \autoref{sec: res_ciclo} relata o teste do ciclo autoral de vantagens, e a \autoref{sec: res_sintese} reúne as respostas. A interpretação desses resultados é feita no \autoref{chap: discussao}.

Salvo indicação em contrário, os valores são medianas das 20 execuções de cada braço, que é a estatística dos testes, e as métricas de equilíbrio vêm da reavaliação independente, com 1000 lutas por par.

\section{Pontos de partida}
\label{sec: res_partida}

O elenco canônico é tão desequilibrado quanto o projeto pretendia. A \autoref{fig: confrontos}(a) mostra a sua matriz de confrontos. O \textit{Rushdown} vence praticamente todas as lutas contra os quatro adversários, e o \textit{Turtle} não vence nenhuma. Nove dos dez pares são \textit{hard counters}, todos decididos a favor de um dos lados em mais de 99\% das lutas, e $\Pdom = \canonDom$, no nível de um elenco sorteado ao acaso. A única aresta moderada é a do \textit{Zoner} contra o \textit{Grappler}, com 63\%. A identidade, por construção, é máxima: desvio nulo, 23 de 23 asserções do validador e concordância $\tau = 1$.

\begin{figure}[htb]
    \caption[Matriz de confrontos do elenco canônico e do algoritmo escalar]{Matriz de confrontos: taxa de vitória, em porcentagem, do personagem da linha contra o da coluna, medida com 16\,000 lutas por par. Em (b), a média das 20 execuções do algoritmo escalar.}
    \label{fig: confrontos}
    \centering
    \includegraphics[width=\textwidth]{confrontos.pdf}
    \source
\end{figure}

A \autoref{tab: referencias} traz os valores das métricas no elenco canônico e nos modelos nulos. Três leituras dela orientam o restante do capítulo. A primeira é que o acaso já produz identidade aparente: um elenco sorteado satisfaz em média 5,4 das 18 asserções estruturais e 1,0 das 5 comportamentais, e o melhor dos 30 sorteados em cada régua chega a 9 e a 3. A segunda é que o piso da concordância $\tau$ é de fato zero, mas com dispersão larga: o melhor nulo chega a \nuloTauMelhor{}. A terceira é que nem a simetria perfeita tem $\Pdom$ nulo. Os espelhos, sem identidade alguma, atingem $\Pdom$ médio de \nuloDomEspelho{}, mas a média mistura duas coisas. O espelho do \textit{Zoner} vale \nuloDomEspelhoZoner{}, e \nuloDomEspelhoZonerDecis{} disso vêm do piso de decisividade: as suas lutas são apertadas demais, não desequilibradas. Os outros quatro espelhos ficam entre \nuloDomEspelhosOutros{}. O termo global médio dos cinco, \nuloGlobEspelho{}, é o que o ruído de medição produz num elenco perfeitamente simétrico, e é contra ele, e não contra zero, que o equilíbrio global de um elenco evoluído deve ser lido.

\begin{table}[htb]
    \caption[Métricas do elenco canônico e dos modelos nulos]{Métricas do elenco canônico e dos modelos nulos (1000 lutas por par para $\Pdom$; 200 para as métricas comportamentais). Na última linha, o melhor dos 30 elencos aleatórios em cada coluna, que não é necessariamente o mesmo elenco.}
    \label{tab: referencias}
    \centering
    \small
    % Fonte: results/baselines/baselines.json
    \begin{tabular}{l c c c c c c}
        \toprule
        \textbf{Elenco} & $\Pdom$ & $\Pdrift$ & \textbf{Cam. 1-2} & \textbf{Cam. 3} & $\tau$ & \textbf{Diferenciação} \\
        & & & (de 18) & (de 5) & & \\
        \midrule
        Canônico              & 1,28  & 0     & 18  & 5   & 1       & 1,35 \\
        Espelhos (média)      & 0,033 & 0,377 & 2,6 & 0,6 & $-$0,18 & 0    \\
        Aleatórios (média)    & 1,33  & 0,415 & 5,4 & 1,0 & $-$0,02 & 1,20 \\
        Aleatórios (melhor)   & 1,05  & 0,327 & 9   & 3   & 0,32    & --   \\
        \bottomrule
    \end{tabular}
    \source
\end{table}

\section{Equilíbrio}
\label{sec: res_equilibrio}

\subsection{Convergência}
\label{sub: res_convergencia}

O algoritmo escalar convergiu em todas as 20 execuções, na geração \convGerMediaAG{} $\pm$ \convGerDpAG{} em média, e sempre entre as gerações \convGerMinAG{} e \convGerMaxAG{}. A \autoref{fig: convergencia} mostra a trajetória do melhor elenco de cada geração. O desequilíbrio cai mais de uma ordem de grandeza nas primeiras 30 gerações, de \trajDomIni{} para \trajDomTrinta{} na mediana, e depois oscila em torno de um patamar, como esperado sob a rotação do \textit{stream}: cada geração é avaliada sob sorteios diferentes. O desvio do melhor elenco cai de cerca de 0,4, o valor de um elenco sorteado, para \trajDriftTrinta{} na geração 30, e depois declina devagar até \trajDriftFim{}.

\begin{figure}[htb]
    \caption[Trajetória do melhor elenco de cada geração]{Trajetória do melhor elenco de cada geração, avaliado dentro do laço. Linhas: mediana das 20 execuções; faixas: intervalo interquartil. O híbrido aparece a partir da geração 75, quando começa a sua fase escalar.}
    \label{fig: convergencia}
    \centering
    \includegraphics[width=\textwidth]{convergencia.pdf}
    \source
\end{figure}

A confirmação independente foi decisiva. Das \gateDisparosAG{} vezes em que o melhor elenco satisfez o predicado de equilíbrio dentro do laço, \gateRecusasAG{} não sobreviveram à reavaliação num \textit{stream} inédito. Convergir, portanto, marca o primeiro sucesso de um teste repetido a cada geração, e não um estado estável. Ao fim do orçamento, \equilAG{} das 20 execuções terminaram com o elenco equilibrado na reavaliação.

O mesmo fenômeno aparece como degradação entre a avaliação no laço e a reavaliação. O $\Pdom$ do elenco final de cada execução do algoritmo escalar tem mediana \domLacoAG{} no laço e \domAG{} na reavaliação, e a mediana da razão entre os dois é \sobreAG{}. A razão acompanha o peso que a seleção dá ao termo ruidoso. Ela vai de \sobreLzero{} no controle $\lambda_{\mathrm{drift}} = 0$, que só otimiza $\Pdom$, a \sobreNSGA{} no representante \textit{scalar\_optimum} do NSGA-II, em que o desvio é um objetivo separado. Entre os braços que otimizam a mesma soma, ela também varia: \sobreAG{} no algoritmo escalar, \sobreSemSem{} no controle sem semente e \sobreHib{} no híbrido, cuja fase final é escalar. A relação com o peso do termo ruidoso é, portanto, uma tendência, e não uma ordenação exata.

\subsection{Equilíbrio reavaliado}
\label{sub: res_equilibrio_reavaliado}

A \autoref{tab: equilibrio} resume o equilíbrio dos elencos finais de cada braço. Todos os personagens de todas as execuções, em todos os braços, terminaram com taxa de vitória global entre 40\% e 60\%. Equilibrar os personagens globalmente, portanto, não discrimina os métodos. O que discrimina são os pares.

\begin{table}[htb]
    \caption[Equilíbrio dos elencos finais]{Equilíbrio dos elencos finais, reavaliados com 1000 lutas por par. $\Pdom$: mediana e, entre colchetes, primeiro e terceiro quartis; $\tglob$: mediana; $\tcap$ e \textit{hard counters}: média por execução.}
    \label{tab: equilibrio}
    \centering
    \footnotesize
    \setlength{\tabcolsep}{4pt}
    % Fonte: results/multi_run/multi_run_{ga,nsga2}.json e results/controls/multi_run_*.json
    \begin{tabular}{l c c c c c}
        \toprule
        \textbf{Braço} & $\Pdom$ & $\tglob$ & $\tcap$ & \textbf{\textit{Hard counters}} & \textbf{Equilibrados} \\
        \midrule
        AG escalar                      & 0,0355 [0,0252; 0,0418] & 0,0318 & 0,0032 & 0,20 & 16/20 \\
        Híbrido                         & 0,0287 [0,0229; 0,0437] & 0,0287 & 0,0042 & 0,25 & 16/20 \\
        NSGA-II (\textit{scalar\_opt.}) & 0,0759 [0,0492; 0,1031] & 0,0385 & 0,0736 & 2,80 & 0/20 \\
        NSGA-II (\textit{best\_dom.}) & 0,0551 [0,0402; 0,0708] & 0,0378 & 0,0467 & 2,05 & 4/20 \\
        Controle $\lambda_{\mathrm{drift}} = 0$ & 0,0432 [0,0303; 0,0524] & 0,0432 & 0,0000 & 0,05 & 19/20 \\
        Controle sem semente            & 0,0264 [0,0183; 0,0481] & 0,0264 & 0,0016 & 0,10 & 18/20 \\
        \bottomrule
    \end{tabular}
    \source
\end{table}

O algoritmo escalar produz elencos cujo $\Pdom$ mediano, \domAG{}, é da ordem do dos espelhos (\nuloDomEspelho{}), o que o põe na posição de 100\% da \autoref{eq: posicao}, entre o piso dos elencos sem estrutura e a simetria perfeita. Essa posição, porém, é dominada pela distância até o piso (\nuloDomPiso{}) e não distingue o elenco evoluído da simetria. Comparados pelo termo global, que não carrega o piso de decisividade do espelho do \textit{Zoner}, os dois se separam: \globAG{} no algoritmo escalar contra \nuloGlobEspelho{} nos espelhos. Em taxa de vitória, os personagens evoluídos se afastam de 50\% em \desvGlobAGpp{} ponto percentual, na raiz da média dos quadrados, contra \desvGlobEspelhopp{} de puro ruído nos espelhos. O algoritmo escalar chega perto da simetria perfeita, mas não a alcança. A \autoref{fig: confrontos}(b) mostra como esse equilíbrio se distribui: na média das execuções, todos os pares ficam entre 46\% e 54\%. Em execuções individuais, dois pares ultrapassaram a tolerância de \textit{hard counter}, cada um em duas das 20 execuções: \textit{Zoner} contra \textit{Turtle} e \textit{Rushdown} contra \textit{Combo Master}.

O desequilíbrio que resta no NSGA-II é de natureza diferente. O seu termo global é próximo do algoritmo escalar (0,0385 contra 0,0318), mas o termo de \textit{hard counters} é vinte vezes maior (0,0736 contra 0,0032). Os dois equilibram os personagens globalmente de modo parecido; o NSGA-II deixa pares ultrapassarem a tolerância, em média \hcNSGAmedia{} por elenco, e nenhuma das suas 20 execuções termina com o elenco equilibrado.

\subsection{Fora das condições de treino}
\label{sub: res_validacao_externa}

A \autoref{tab: validacao_externa} traz a validação externa dos quatro elencos de semente 42. O elenco do algoritmo escalar é o único que replica: sob sorteios inéditos e 5000 lutas por par, todas as 15 taxas de vitória ficam dentro das faixas de equilíbrio, com $\Pdom = \extDomReplAG{}$, contra \extDomLacoAG{} dentro do laço. O representante \textit{scalar\_optimum} do NSGA-II não replica: o par \textit{Rushdown} contra \textit{Combo Master} fica em 32,7\%, com o intervalo de confiança inteiro fora da faixa. O \textit{knee\_point} e o canônico falham com larga margem.

\begin{table}[htb]
    \caption[Validação externa dos elencos de semente 42]{Validação externa dos elencos de semente 42: veredito (R: robusto; I: inconclusivo; F: frágil) e $\Pdom$ sob cada condição, ambos calculados sobre 5000 lutas por par (10 sementes de 500 lutas, somadas).}
    \label{tab: validacao_externa}
    \centering
    \footnotesize
    \setlength{\tabcolsep}{4pt}
    % Fonte: results/external_validation/external_validation_{evolved,nsga2_scalar_optimum,nsga2_knee_point,canonical}.json
    \begin{tabular}{l c c c c}
        \toprule
        \textbf{Condição} & \textbf{AG escalar} & \textbf{NSGA-II, \textit{scalar\_opt.}} & \textbf{NSGA-II, \textit{knee}} & \textbf{Canônico} \\
        \midrule
        Replicação (regras do treino) & R \; 0,033 & F \; 0,051 & F \; 0,400 & F \; 1,277 \\
        \midrule
        Distância inicial 40        & R \; 0,025 & F \; 0,068 & F \; 0,394 & F \; 1,277 \\
        Distância inicial 60        & R \; 0,040 & F \; 0,045 & F \; 0,396 & F \; 1,278 \\
        Campo 80                    & F \; 0,087 & F \; 0,143 & F \; 0,470 & F \; 1,297 \\
        Campo 120                   & R \; 0,033 & F \; 0,160 & F \; 0,407 & F \; 1,272 \\
        Persistência 4              & I \; 0,110 & F \; 0,273 & F \; 0,440 & F \; 1,283 \\
        Persistência 6              & F \; 0,153 & F \; 0,230 & F \; 0,441 & F \; 1,271 \\
        Multiplicador da guarda 0,55 & I \; 0,075 & F \; 0,108 & F \; 0,397 & F \; 1,277 \\
        Multiplicador da guarda 0,65 & R \; 0,061 & F \; 0,063 & F \; 0,444 & F \; 1,278 \\
        \bottomrule
    \end{tabular}
    \source
\end{table}

Sob regras que a busca nunca viu, o elenco do algoritmo escalar é robusto em quatro das oito condições, inconclusivo em duas e frágil em duas. As falhas se concentram nos pares próximos da borda da faixa. No campo de 80 unidades, o \textit{Zoner} passa a vencer o \textit{Turtle} em 68,8\% das lutas; com persistência 6, o \textit{Combo Master} vence o \textit{Turtle} em 69,4\% e a taxa global do \textit{Turtle} cai a 39,3\%. Nas duas condições inconclusivas, um único par tem o intervalo de confiança cruzando a borda de 35\%. A persistência da intenção é a regra à qual o equilíbrio é mais sensível: as duas perturbações dela produzem os maiores $\Pdom$ da coluna.

\subsection{O que a busca enxerga}
\label{sub: res_sensibilidade}

A \autoref{fig: sensibilidade} mostra a análise de sensibilidade sobre o elenco do algoritmo escalar de semente 42. O alcance, o dano, o intervalo entre golpes e os pontos de vida movem a taxa de vitória com força: deslocar o alcance numa janela de um desvio de mutação para cada lado muda a taxa global dos personagens em 30,8 pontos percentuais em média, contra um piso de ruído medido de \sensPiso{} pontos. O agarrão e o atordoamento também são claramente visíveis. O peso de RECUAR, a repulsão e a velocidade ficam na faixa limítrofe, e os pesos de FRENTE e de GUARDA, no fundo da ordenação, não se distinguem do ruído. A janela dos pesos comportamentais, porém, é quatro vezes mais estreita que a dos atributos, porque o desvio da mutação deles é menor (\autoref{sec: ag_escalar}).

\begin{figure}[htb]
    \caption[Sensibilidade de cada gene no elenco do algoritmo escalar]{Sensibilidade de cada gene no elenco do algoritmo escalar de semente 42: variação média absoluta da taxa de vitória global dos cinco personagens numa janela de dois desvios-padrão de mutação.}
    \label{fig: sensibilidade}
    \centering
    \includegraphics[width=\textwidth]{sensibilidade.pdf}
    \source
\end{figure}

A diferença de janela separa duas explicações para o fundo da ordenação: a política pesar pouco no combate, ou a mutação mexer pouco nela. Numa medição feita para este texto, os pesos foram deslocados com a mesma janela relativa dos atributos, de 10\% da amplitude para cada lado. Nessa condição, os três pesos movem a taxa global em \sensPesosIguais{} pontos percentuais (\sensWRIgual{} no de RECUAR, \sensWFIgual{} no de FRENTE e \sensWGIgual{} no de GUARDA), acima do atordoamento (\sensStun{}), da repulsão (\sensKb{}) e da velocidade (\sensSpd{}). A política pesa no combate tanto quanto os atributos intermediários. O que a esconde da busca é o passo com que a mutação a altera: nesse passo, o efeito de cada mutação sobre o equilíbrio fica no nível do ruído, e o único termo da aptidão que ainda distingue uma mutação de outra é o de identidade. Essa assimetria reaparece na \autoref{sub: res_perdas}.

\section{Identidade}
\label{sec: res_identidade}

\subsection{Identidade do algoritmo escalar}
\label{sub: res_identidade_ag}

A \autoref{fig: metricas} mostra a distribuição, nas 20 execuções de cada braço, das seis métricas da família de testes. Tomado isoladamente, o algoritmo escalar preserva a identidade estrutural com folga sobre o acaso: desvio mediano de \driftAG{}, contra 0,415 dos elencos sorteados, e \LumdoisAG{} das 18 asserções das Camadas 1 e 2, contra 5,4 em média e 9 no melhor nulo. Essa régua, contudo, é parcialmente endógena, e superar os nulos nela é esperado de qualquer busca com o termo de desvio.

\begin{figure}[htbp]
    \caption[As seis métricas da família de testes por braço]{As seis métricas da família de testes nas 20 execuções de cada braço. Pontos: execuções; barra preta: mediana. Linha tracejada: média dos espelhos em (a) e dos elencos aleatórios de (c) a (f). O NSGA-II aparece pelo representante \textit{scalar\_optimum}.}
    \label{fig: metricas}
    \centering
    \includegraphics[width=\textwidth]{metricas.pdf}
    \source
\end{figure}

A identidade funcional é mais modesta. A mediana é de \LtresAG{} das 5 asserções comportamentais e $\tau = \tauAG{}$. Os dois valores estão acima da média dos elencos sorteados (1,0 e $-0{,}02$), mas no topo da distribuição dos nulos: o melhor elenco sorteado também atinge 3 asserções comportamentais e $\tau = \nuloTauMelhor{}$. Um elenco evoluído isolado, comparado a 35 nulos, não tem resolução para concluir. O elenco de semente 42, por exemplo, empata com o melhor nulo na Camada 3 ($p = 0{,}03$) e fica um fio abaixo dele na concordância ($\tau = 0{,}314$, $p = 0{,}06$). A evidência sobre a identidade funcional vem, por isso, da comparação com o controle.

\subsection{O contrafactual: equilibrar sem o termo de identidade}
\label{sub: res_controle}

A \autoref{tab: controles} compara o algoritmo escalar com os dois braços de controle. Sem o termo de identidade, o mesmo algoritmo, com a mesma amostra e o mesmo orçamento, perde a identidade nas quatro réguas, todas com efeito grande e $p_{\mathrm{Holm}} \le \pHolmLzeroIdentMax{}$. O desvio mediano sobe de \driftAG{} para \driftLzero{}; as asserções estruturais caem de \LumdoisAG{} para \LumdoisLzero{}; as comportamentais, de \LtresAG{} para \LtresLzero{}; e a concordância cai de \tauAG{} para \tauLzero{}, com média de \tauLzeroMedia{} nas 20 execuções.

\begin{table}[htb]
    \caption[Algoritmo escalar contra os dois controles]{Algoritmo escalar contra os dois controles: medianas, $p$ corrigido por Holm (Wilcoxon pareado) e $\hat{A}_{12}$ (probabilidade de o AG escalar dar valor maior).}
    \label{tab: controles}
    \centering
    \small
    % Fonte: results/controls/comparison_ga_vs_ga_drift0_dom1.json e comparison_ga_vs_ga_unseeded.json
    \begin{tabular}{l c c c c c c c}
        \toprule
        & & \multicolumn{3}{c}{\textbf{Controle $\lambda_{\mathrm{drift}} = 0$}} & \multicolumn{3}{c}{\textbf{Controle sem semente}} \\
        \cmidrule(lr){3-5} \cmidrule(lr){6-8}
        \textbf{Métrica} & \textbf{AG} & \textbf{Med.} & $p_{\mathrm{Holm}}$ & $\hat{A}_{12}$ & \textbf{Med.} & $p_{\mathrm{Holm}}$ & $\hat{A}_{12}$ \\
        \midrule
        $\Pdom$              & 0,0355 & 0,0432  & 0,059                & 0,30 & 0,0264 & 0,99   & 0,55 \\
        \textit{Hard counters} & 0    & 0       & 0,18                 & 0,57 & 0      & 0,83   & 0,55 \\
        $\Pdrift$            & 0,2473 & 0,4055  & $1{,}1\cdot10^{-5}$  & 0,00 & 0,2703 & 0,0073 & 0,25 \\
        Camadas 1-2          & 11     & 6       & $4{,}2\cdot10^{-4}$  & 0,98 & 10,5   & 0,31   & 0,66 \\
        Camada 3             & 3      & 1       & 0,0044               & 0,85 & 2      & 0,59   & 0,64 \\
        $\tau$               & 0,28   & $-$0,03 & $6{,}7\cdot10^{-4}$  & 0,87 & 0,19   & 0,39   & 0,67 \\
        \bottomrule
    \end{tabular}
    \source
\end{table}

O controle não perde apenas \textit{parte} da identidade: ele cai ao nível do acaso. O seu desvio (\driftLzero{}) é o de um elenco sorteado (0,415); as suas asserções estruturais (6) e comportamentais (1) estão na média dos elencos sorteados; a sua concordância média, \tauLzeroMedia{}, é a do acaso. A \autoref{fig: metricas} mostra o mesmo em cada execução: nas quatro réguas de identidade, a nuvem do controle fica sobre a linha dos elencos aleatórios. Sem o termo de identidade, o equilíbrio é alcançado com elencos que não guardam relação com os arquétipos.

No equilíbrio, os dois braços não se separam. O $\Pdom$ mediano do controle é maior (\domLzero{} contra \domAG{}), com efeito médio ($\hat{A}_{12} = \aLzeroDom{}$) e $p$ bruto de \pBrutoLzeroDom{}, mas a diferença não sobrevive à correção para múltiplas comparações ($p_{\mathrm{Holm}} = \pHolmLzeroDom{}$). O número de \textit{hard counters} também não separa. E, descritivamente, o controle termina equilibrado em \equilLzero{} das 20 execuções, contra \equilAG{} do algoritmo escalar. O resultado sustenta que a identidade não custou equilíbrio mensurável. Não sustenta que ela o melhore.

O controle sem a semente canônica separa do algoritmo escalar apenas no desvio (\driftSemSem{} contra \driftAG{}, $p_{\mathrm{Holm}} = \pHolmSemSemDrift{}$), e em nenhuma outra métrica. Nas demais réguas de identidade, as estimativas também favorecem a semente, com efeito médio nas Camadas 1 e 2 e na concordância $\tau$ (\autoref{tab: controles}), mas sem significância. A semente contribui um pouco para a identidade e não explica as diferenças entre o algoritmo escalar e o NSGA-II relatadas na \autoref{sec: res_compromisso}.

\subsection{O que se preserva e o que se perde}
\label{sub: res_perdas}

A \autoref{tab: assercoes} desagrega o validador: em quantas das 20 execuções de cada braço cada asserção foi satisfeita. Três padrões aparecem.

\begin{table}[htbp]
    \caption[Asserções do validador satisfeitas por braço]{Número de execuções, em 20, em que cada asserção do validador foi satisfeita. Para as asserções das Camadas 1 e 3, o acaso dá em torno de 4 em 20. O NSGA-II aparece pelo representante \textit{scalar\_optimum}.}
    \label{tab: assercoes}
    \centering
    \footnotesize
    % Fonte: validador reaplicado aos genes gravados em results/multi_run e results/controls
    % (200 lutas por par, semente 9999); reproduz os totais gravados nas 80 execuções.
    \begin{tabularx}{\textwidth}{l X c c c c}
        \toprule
        \textbf{Arquétipo} & \textbf{Asserção} & \textbf{AG} & \textbf{Híbrido} & \textbf{NSGA-II} & $\boldsymbol{\lambda = 0}$ \\
        \midrule
        \multicolumn{6}{l}{\textit{Camada 1: genes, entre personagens}} \\
        \textit{Rushdown}     & maior velocidade                     & 16 & 19 & 17 & 3 \\
                              & menor intervalo entre golpes         & 16 & 19 & 19 & 8 \\
                              & maior probabilidade de FRENTE        & 5  & 4  & 5  & 4 \\
        \textit{Zoner}        & maior alcance                        & 11 & 16 & 17 & 4 \\
                              & maior repulsão                       & 12 & 14 & 18 & 5 \\
                              & maior probabilidade de RECUAR        & 9  & 15 & 17 & 6 \\
        \textit{Combo Master} & maior atordoamento                   & 16 & 18 & 18 & 2 \\
        \textit{Grappler}     & maior dano                           & 10 & 19 & 19 & 4 \\
                              & maior poder de agarrão               & 16 & 18 & 19 & 4 \\
        \textit{Turtle}       & menor velocidade                     & 12 & 16 & 17 & 5 \\
                              & maior intervalo entre golpes         & 9  & 13 & 15 & 4 \\
                              & mais pontos de vida                  & 10 & 19 & 20 & 4 \\
                              & maior probabilidade de GUARDA        & 10 & 15 & 16 & 2 \\
        \midrule
        \multicolumn{6}{l}{\textit{Camada 2: genes, no mesmo personagem}} \\
        \textit{Zoner}        & alcance maior que velocidade         & 15 & 19 & 19 & 8 \\
        \textit{Rushdown}     & velocidade maior que alcance         & 20 & 20 & 20 & 16 \\
        \textit{Grappler}     & pontos de vida maiores que velocidade & 11 & 18 & 19 & 7 \\
        \textit{Combo Master} & atordoamento maior que repulsão      & 18 & 20 & 20 & 13 \\
                              & velocidade maior que alcance         & 14 & 13 & 14 & 15 \\
        \midrule
        \multicolumn{6}{l}{\textit{Camada 3: comportamento}} \\
        \textit{Zoner}        & maior distância média de luta        & 10 & 16 & 17 & 2 \\
        \textit{Rushdown}     & mais golpes conectados               & 3  & 10 & 11 & 6 \\
        \textit{Turtle}       & maior fração de defesa escolhida     & 11 & 15 & 16 & 2 \\
        \textit{Combo Master} & maior atordoamento infligido         & 12 & 16 & 14 & 3 \\
        \textit{Grappler}     & maior dano através da guarda         & 13 & 17 & 15 & 5 \\
        \bottomrule
    \end{tabularx}
    \source
\end{table}

O primeiro padrão é que o controle $\lambda_{\mathrm{drift}} = 0$ fica no nível do acaso asserção por asserção, e não apenas no total. O segundo é que o algoritmo escalar preserva melhor os traços de ritmo do \textit{Rushdown} (velocidade e intervalo entre golpes), o atordoamento do \textit{Combo Master} e o agarrão do \textit{Grappler}, todos em 16 de 20 execuções, e perde mais os traços do \textit{Zoner} e do \textit{Turtle}, que ficam entre 9 e 12. O terceiro, e o mais nítido, é que a asserção da política do \textit{Rushdown}, ser o personagem de maior probabilidade de avançar, é satisfeita em apenas 4 a 5 de 20 execuções em \textit{todos} os braços, inclusive nos que melhor preservam a identidade. Em boa parte dessas falhas, o \textit{Combo Master}, cuja probabilidade canônica de FRENTE (0,74) já é próxima da do \textit{Rushdown}, passa a avançar mais que ele. Isso ocorre em \cmAvancaMaisAG{} das 20 execuções do algoritmo escalar, em \cmAvancaMaisNSGA{} das do NSGA-II e em \cmAvancaMaisHib{} das do híbrido. A asserção comportamental correspondente, ser o que mais golpeia, cai para 3 em 20 no algoritmo escalar.

A \autoref{tab: politica} mostra a origem desse padrão. Ela compara, para cada arquétipo, a probabilidade da intenção que o define no elenco canônico com a mediana dessa probabilidade nas execuções de cada braço. Nos elencos evoluídos, a política de todos os personagens se afasta do perfil canônico, mas não em direção a uma política sorteada ao acaso, que daria 1/3 a cada intenção. Os personagens se aproximam de um perfil comum, que recua pouco: no algoritmo escalar, os quatro arquétipos além do \textit{Zoner} recuam em \polAGRecuoOutros{} das intenções. O \textit{Rushdown} canônico avança em 86\% das intenções; o do algoritmo escalar, em 39\%, e o do NSGA-II, em 48\%. No controle sem termo de identidade, as políticas dos cinco arquétipos são praticamente indistinguíveis entre si: na mediana, todos avançam em \polLzeroF{} das intenções, recuam em \polLzeroR{} e guardam em \polLzeroG{}. O valor baixo do \textit{Zoner} na \autoref{tab: politica} não o distingue: é a propensão a recuar comum aos cinco.

\begin{table}[htb]
    \caption[Probabilidade da intenção que define cada arquétipo]{Probabilidade da intenção que define cada arquétipo: valor canônico e mediana das 20 execuções de cada braço. Uma política sorteada ao acaso dá, em média, 1/3 a cada intenção. O NSGA-II aparece pelo representante \textit{scalar\_optimum}.}
    \label{tab: politica}
    \centering
    \small
    % Fonte: genes gravados em results/multi_run e results/controls (probabilidades da Eq. de intenção)
    \begin{tabular}{l l c c c c c}
        \toprule
        \textbf{Arquétipo} & \textbf{Intenção} & \textbf{Canônico} & \textbf{AG} & \textbf{Híbrido} & \textbf{NSGA-II} & $\boldsymbol{\lambda = 0}$ \\
        \midrule
        \textit{Rushdown}     & FRENTE & 0,86 & 0,39 & 0,47 & 0,48 & 0,35 \\
        \textit{Combo Master} & FRENTE & 0,74 & 0,47 & 0,51 & 0,55 & 0,40 \\
        \textit{Grappler}     & FRENTE & 0,58 & 0,41 & 0,44 & 0,44 & 0,36 \\
        \textit{Zoner}        & RECUAR & 0,55 & 0,33 & 0,45 & 0,46 & 0,19 \\
        \textit{Turtle}       & GUARDA & 0,54 & 0,45 & 0,47 & 0,48 & 0,40 \\
        \bottomrule
    \end{tabular}
    \source
\end{table}

A diferenciação, calculada para este texto a partir dos genes gravados, não acrescenta evidência. Ela cai de \difCanon{} no canônico para \difAG{} na mediana do algoritmo escalar, abaixo da média dos elencos sorteados (\difAleat{}), enquanto o híbrido (\difHib{}) e o NSGA-II (\difNSGA{}) ficam no nível dos sorteados. O controle sem termo de identidade, porém, fica em \difLzero{}, acima do algoritmo escalar. Se o equilíbrio fosse o que aproxima os personagens, o braço que só otimiza equilíbrio teria o menor valor. Como elencos sorteados são dispersos sem ter identidade alguma, a métrica não separa identidade de dispersão, e não serve como evidência de homogeneização pelo equilíbrio.

\section{O compromisso entre equilíbrio e identidade}
\label{sec: res_compromisso}

\subsection{A fronteira de Pareto}
\label{sub: res_fronteira}

A \autoref{fig: fronteira} sobrepõe as 20 fronteiras finais do NSGA-II e os elencos finais dos braços escalares, todos avaliados dentro do laço, sob o \textit{stream} da última geração de cada semente. As fronteiras têm em média \fronteiraTamMedio{} pontos e são notavelmente estáveis entre execuções: o hipervolume é \hvMedia{} $\pm$ \hvDp{} e o espaçamento, \espacMedia{} $\pm$ \espacDp{}. O formato é o de um compromisso íngreme. Na ponta de identidade, os elencos mantêm desvio próximo de 0,06, mas com $\Pdom$ acima de 1, tão desequilibrados quanto o canônico. Para reduzir $\Pdom$ a um décimo desse valor, a fronteira paga de 0,08 a 0,1 em desvio.

\begin{figure}[htb]
    \caption[Fronteiras finais do NSGA-II e elencos finais dos demais braços]{Fronteiras finais do NSGA-II (20 execuções) e elencos finais dos demais braços, todos avaliados dentro do laço, sob o \textit{stream} da última geração. O elenco canônico aparece pela sua reavaliação com 1000 lutas por par.}
    \label{fig: fronteira}
    \centering
    \includegraphics[width=\textwidth]{fronteira.pdf}
    \source
\end{figure}

A \autoref{tab: representantes} resume os cinco representantes extraídos de cada fronteira. Os três pontos geométricos, o extremo de identidade, o joelho e o ideal, preservam quase toda a identidade (Camada 3 em 5 de 5, $\tau$ de 0,70 a 0,80), mas nenhum deles termina equilibrado em nenhuma execução. Os pontos de interesse para a pergunta de pesquisa estão na ponta de equilíbrio da fronteira.

\begin{table}[htb]
    \caption[Representantes das fronteiras do NSGA-II]{Representantes das fronteiras do NSGA-II, reavaliados: medianas das 20 execuções; \textit{hard counters} em média por execução.}
    \label{tab: representantes}
    \centering
    \small
    % Fonte: results/multi_run/multi_run_nsga2.json (representatives)
    \begin{tabular}{l c c c c c c c}
        \toprule
        \textbf{Representante} & $\Pdom$ & $\Pdrift$ & \textbf{\textit{H. c.}} & \textbf{Equil.} & \textbf{Cam. 1-2} & \textbf{Cam. 3} & $\tau$ \\
        \midrule
        \textit{best\_drift}     & 1,056 & 0,059 & 8,45 & 0/20 & 18   & 5 & 0,80 \\
        \textit{knee\_point}     & 0,405 & 0,087 & 6,85 & 0/20 & 17   & 5 & 0,70 \\
        \textit{ideal\_point}    & 0,396 & 0,087 & 7,25 & 0/20 & 17,5 & 5 & 0,71 \\
        \textit{scalar\_optimum} & 0,076 & 0,145 & 2,80 & 0/20 & 15   & 4 & 0,54 \\
        \textit{best\_dominance} & 0,055 & 0,166 & 2,05 & 4/20 & 15   & 4 & 0,51 \\
        \bottomrule
    \end{tabular}
    \source
\end{table}

\subsection{Algoritmo escalar e NSGA-II}
\label{sub: res_ag_nsga}

A \autoref{tab: comparacoes} traz a comparação estatística do algoritmo escalar com o NSGA-II e com o híbrido. Contra o \textit{scalar\_optimum}, as seis métricas da família são significativas, todas com efeito grande ($p_{\mathrm{Holm}} \le \pHolmNSGAmax{}$), e a divisão é limpa. O algoritmo escalar vence as duas métricas de equilíbrio: $\Pdom$ de \domAG{} contra \domNSGA{} e nenhum \textit{hard counter} na mediana, contra 3. O NSGA-II vence as quatro de identidade: desvio de \driftNSGA{} contra \driftAG{}, com separação completa ($\hat{A}_{12} = 1{,}00$), \LumdoisNSGA{} asserções estruturais contra \LumdoisAG{}, \LtresNSGA{} comportamentais contra \LtresAG{} e concordância de \tauNSGA{} contra \tauAG{}. Contra o extremo \textit{best\_dominance}, as seis diferenças mantêm a direção e a significância, com efeitos um pouco menores em equilíbrio. Em todas as comparações, o número de personagens em banda foi constante (5 em todas as execuções) e ficou fora da família de testes, que teve seis métricas.

\begin{table}[htb]
    \caption[Algoritmo escalar contra o NSGA-II e o híbrido]{Algoritmo escalar contra o NSGA-II e o híbrido: medianas, $p$ corrigido por Holm (Wilcoxon pareado) e $\hat{A}_{12}$ (probabilidade de o AG escalar dar valor maior). Todos os efeitos com $p_{\mathrm{Holm}} < 0{,}05$ são grandes.}
    \label{tab: comparacoes}
    \centering
    \footnotesize
    \setlength{\tabcolsep}{4pt}
    % Fonte: results/multi_run/comparison_ga_vs_nsga2*.json e results/controls/comparison_ga_vs_hybrid_split0.5_front.json
    \begin{tabular}{l c ccc ccc ccc}
        \toprule
        & & \multicolumn{3}{c}{\textbf{NSGA-II, \textit{scalar\_opt.}}} & \multicolumn{3}{c}{\textbf{NSGA-II, \textit{best\_dom.}}} & \multicolumn{3}{c}{\textbf{Híbrido}} \\
        \cmidrule(lr){3-5} \cmidrule(lr){6-8} \cmidrule(lr){9-11}
        \textbf{Métrica} & \textbf{AG} & \textbf{Med.} & $p_{\mathrm{Holm}}$ & $\hat{A}_{12}$ & \textbf{Med.} & $p_{\mathrm{Holm}}$ & $\hat{A}_{12}$ & \textbf{Med.} & $p_{\mathrm{Holm}}$ & $\hat{A}_{12}$ \\
        \midrule
        $\Pdom$                & 0,0355 & 0,0759 & $9{,}5\cdot10^{-5}$ & 0,07 & 0,0551 & 0,0034              & 0,18 & 0,0287 & 1,00                & 0,51 \\
        \textit{Hard counters} & 0      & 3      & $3{,}4\cdot10^{-4}$ & 0,03 & 2      & 0,0021              & 0,14 & 0      & 1,00                & 0,49 \\
        $\Pdrift$              & 0,2473 & 0,1445 & $1{,}1\cdot10^{-5}$ & 1,00 & 0,1663 & $2{,}3\cdot10^{-5}$ & 0,98 & 0,1663 & $1{,}1\cdot10^{-5}$ & 0,97 \\
        Camadas 1-2            & 11     & 15     & $3{,}4\cdot10^{-4}$ & 0,05 & 15     & 0,0013              & 0,09 & 15     & 0,0026              & 0,12 \\
        Camada 3               & 3      & 4      & 0,0027              & 0,24 & 4      & 0,0034              & 0,24 & 4      & 0,0058              & 0,24 \\
        $\tau$                 & 0,28   & 0,54   & $1{,}1\cdot10^{-4}$ & 0,09 & 0,51   & $5{,}2\cdot10^{-4}$ & 0,13 & 0,52   & $8{,}4\cdot10^{-4}$ & 0,15 \\
        \bottomrule
    \end{tabular}
    \source
\end{table}

Nenhuma dessas conclusões depende da escolha do teste. O teste não pareado de Mann-Whitney, reportado ao lado nos artefatos, dá $p$ brutos menores que 0,004 nas seis métricas da comparação com o \textit{scalar\_optimum}.

\subsection{Relação de Pareto por semente}
\label{sub: res_pareto}

A comparação por representante mede, em parte, a escolha do representante. A relação de Pareto por semente não depende dela. Em 18 das 20 sementes, o ponto do algoritmo escalar e a fronteira do NSGA-II são mutuamente não dominados; em uma, o ponto do escalar domina algum ponto da fronteira; e em outra, é dominado por ela. A \autoref{fig: fronteira} mostra por quê: em 19 das 20 sementes, o ponto do algoritmo escalar está \textit{além} da ponta de equilíbrio da fronteira, com $\Pdom$ no laço de \domLacoAG{} contra \domLacoFronteiraMin{} no ponto mais equilibrado da fronteira. O escalar alcança um equilíbrio que a fronteira não oferece, pagando em desvio.

Ser não dominado, contudo, não é ser ótimo na própria função. Na soma $\Pdom + \Pdrift$, que é exatamente o que o algoritmo escalar maximiza com sinal trocado (\autoref{eq: aptidao_escalar}), a fronteira do NSGA-II contém um ponto melhor que o do escalar em 20 das 20 sementes, com medianas de \somaFronteira{} contra \somaAG{} dentro do laço. A avaliação no laço favorece mais o algoritmo escalar que o NSGA-II (\autoref{sub: res_convergencia}), e a conclusão se mantém com os elencos reavaliados: o representante \textit{scalar\_optimum} tem soma menor que a do escalar em \somaReavSOmelhorAG{} das 20 sementes. O algoritmo escalar termina num extremo do compromisso, e não no ponto que a sua própria função de aptidão preferiria.

\subsection{O híbrido}
\label{sub: res_hibrido}

A última coluna da \autoref{tab: comparacoes} compara o híbrido com o algoritmo escalar. As duas métricas de equilíbrio não se movem: $\hat{A}_{12}$ de 0,51 e 0,49, $p_{\mathrm{Holm}} = 1{,}00$, e \equilHib{} elencos equilibrados em 20, como no escalar. As quatro métricas de identidade melhoram, todas com efeito grande e $p_{\mathrm{Holm}} \le \pHolmHibIdentMax{}$: o desvio cai de \driftAG{} para \driftHib{}; as asserções estruturais sobem de \LumdoisAG{} para \LumdoisHib{} e as comportamentais, de \LtresAG{} para \LtresHib{}; a concordância quase dobra, de \tauAG{} para \tauHib{}.

Descritivamente, o híbrido reúne a metade favorável de cada algoritmo. A sua identidade fica próxima da do NSGA-II (desvio \driftHib{} contra \driftNSGA{}, concordância \tauHib{} contra \tauNSGA{}, Camada 3 igual), com o equilíbrio do algoritmo escalar: 0,25 \textit{hard counters} por elenco contra \hcNSGAmedia{}, e \equilHib{} elencos equilibrados contra \equilNSGA{}. Na \autoref{fig: fronteira}, os pontos do híbrido ficam, como os do algoritmo escalar, à esquerda da ponta de equilíbrio das fronteiras (em \hibAlemPonta{} das 20 sementes), mas com desvio menor que o deles. Em nenhuma semente o ponto do híbrido é dominado pela fronteira do NSGA-II da mesma semente, e em \hibDominaFronteira{} das 20 ele domina algum ponto dela. Ao contrário do algoritmo escalar, o híbrido termina perto do ótimo da própria função. A fronteira contém um ponto de soma $\Pdom + \Pdrift$ menor que a dele em apenas \hibFronteiraMelhor{} das 20 sementes, contra 20 de 20 no escalar, e a mediana da sua soma no laço, \somaHib{}, fica abaixo da do melhor ponto da fronteira, \somaFronteira{}. Com os elencos reavaliados, o híbrido tem soma menor que a do representante \textit{scalar\_optimum} em \somaReavHibMelhorSO{} das 20 sementes. Essa comparação com o NSGA-II é descritiva: a família de testes pré-definida opõe o algoritmo escalar a cada braço, e não os braços entre si.

O preço do híbrido é a velocidade. Ele converge, em média, na geração \convGerMediaHib{} $\pm$ \convGerDpHib{}, contra \convGerMediaAG{} $\pm$ \convGerDpAG{} do escalar, porque as suas primeiras 75 gerações são de Pareto e não perseguem o predicado de equilíbrio. A \autoref{fig: convergencia} mostra a consequência: ao receber a fronteira, a fase escalar parte de um $\Pdom$ mais alto que o do escalar na mesma geração, aproxima-se dele nas gerações seguintes e mantém o desvio no nível do NSGA-II até o fim.

\section{Estrutura de vantagens}
\label{sec: res_ciclo}

A \autoref{tab: ciclo} resume o teste do ciclo autoral, com 16\,000 lutas por par.

\begin{table}[htb]
    \caption[Teste do ciclo autoral de vantagens]{Teste do ciclo autoral de vantagens (16\,000 lutas por par): arestas decididas por elenco, arestas mantidas entre as decididas, tríades circulares e margem mediana das arestas de cada elenco, em média sobre os elencos. O NSGA-II aparece pelo representante \textit{scalar\_optimum}.}
    \label{tab: ciclo}
    \centering
    \small
    % Fonte: results/cycle/cycle_structure.json
    \begin{tabular}{l c c c c c}
        \toprule
        \textbf{Grupo} & \textbf{Elencos} & \textbf{Decididas} & \textbf{Mantidas} & \textbf{Tríades} & \textbf{Margem} \\
        & & (de 10) & (das decididas) & (de 5) & \\
        \midrule
        Canônico    & 1  & 10,00 & 6/10 (60\%)        & 1,00 & 0,500 \\
        AG escalar  & 20 & 9,30  & 105/186 (56,5\%)   & 3,71 & 0,048 \\
        NSGA-II     & 20 & 9,85  & 98/197 (49,7\%)    & 4,17 & 0,104 \\
        Espelhos    & 5  & 1,00  & 3/5                & 2,59 & 0,003 \\
        Aleatórios  & 30 & 10,00 & 149/300 (49,7\%)   & 0,40 & 0,491 \\
        \bottomrule
    \end{tabular}
    \source
\end{table}

O primeiro resultado é sobre o próprio canônico: ele realiza apenas 6 das 10 arestas autorais. As quatro que quebra são inversões completas: o \textit{Combo Master} vence o \textit{Grappler} em 0,2\% das lutas, e não na maioria; o \textit{Grappler} vence o \textit{Rushdown} em 0,6\%; e o \textit{Turtle} não vence nem o \textit{Rushdown} nem o \textit{Combo Master}. Lidas juntas com a \autoref{fig: confrontos}(a), essas arestas descrevem uma hierarquia, e não um ciclo: o \textit{Rushdown} vence todos e o \textit{Turtle} perde para todos. As tríades circulares confirmam, com \triadesCanon{} de um máximo de 5.

Nos elencos evoluídos, o ciclo não é realizado. O algoritmo escalar mantém \cicloMantidasAG{} das \cicloDecididasAG{} arestas decididas, \cicloTaxaAG{}, uma inclinação na direção autoral que não é significativa: $p = \cicloBinomAG{}$ no teste binomial, e $p = \cicloPNulos{}$ ($\hat{A}_{12} = \cicloANulos{}$) contra os elencos aleatórios, que ficam em \cicloTaxaNulos{}. O NSGA-II mantém praticamente metade. A aresta mais emblemática, o agarrão do \textit{Grappler} contra a guarda do \textit{Turtle}, vence 100\% das lutas no canônico e fica em 52\% na média das execuções do algoritmo escalar.

As tríades mostram onde está a estrutura não transitiva. Elas são quase nulas nos elencos aleatórios (\triadesAleat{}), mínimas no canônico (\triadesCanon{}) e altas nos elencos evoluídos: \triadesAG{} no algoritmo escalar e \triadesNSGA{} no NSGA-II, com quase todas as arestas decididas. Os espelhos, com \triadesEspelho{} tríades e apenas uma aresta decidida por elenco, mostram que a contagem só tem significado quando as arestas são decididas.

\section{Síntese}
\label{sec: res_sintese}

O \autoref{quad: sintese} reúne as respostas às três questões de pesquisa, na forma mais estreita que os dados sustentam.

\begin{quadro}[htb]
    \caption{Respostas às questões de pesquisa.}
    \label{quad: sintese}
    \centering
    \small
    \begin{tabularx}{\textwidth}{l X}
        \toprule
        \textbf{Questão} & \textbf{Resposta} \\
        \midrule
        QP1: equilíbrio & Sim. O algoritmo escalar converge em todas as execuções e termina, na mediana, perto da simetria perfeita, com o elenco equilibrado em 16 de 20. O elenco da semente 42, único submetido à validação externa, replica fora do laço e é robusto a 4 de 8 regras não vistas, frágil em 2. \\
        \midrule
        QP2: identidade & Parcialmente, e a parte preservada é atribuível ao termo de identidade. Sem ele, a identidade cai ao nível do acaso nas quatro réguas, sem ganho de equilíbrio significativo. Com ele, a identidade funcional fica acima do acaso, mas longe do canônico, e a política dos personagens é o componente menos preservado. \\
        \midrule
        QP3: compromisso & O algoritmo escalar ocupa a ponta de equilíbrio do compromisso, além da fronteira do NSGA-II; o NSGA-II preserva mais identidade e deixa pares desequilibrados. Nenhum domina o outro. O híbrido, com o mesmo orçamento, aproxima-se da identidade do NSGA-II mantendo o equilíbrio do escalar. \\
        \bottomrule
    \end{tabularx}
    \source
\end{quadro}
